\section{Datasets}

The following chapter contains an overview of the datasets and the observational filters commonly applied across all validation runs in this report.
\\[1em]
\noindent The common period for all datasets is from \textbf{\detokenize{$<ReportVars['Common']['interval_from']>$} to \detokenize{$<ReportVars['Common']['interval_to']>$}}.
Only measurements within this period are used for validation.

\subsection{Candidate - \detokenize{$<Run1ContentVars['ConfigVars']['DS1']['name']>$} (\detokenize{$<Run1ContentVars['ConfigVars']['DS1']['version']>$}) Soil Moisture}

\noindent The SMOS Level 2 (L2) Soil Moisture product is a global dataset of surface
soil moisture (and related vegetation parameters) derived from ESA's SMOS satellite
\citep{kerr2001}.
SMOS carries a passive L-band microwave radiometer (MIRAS) that measures brightness
temperatures emitted from the Earth's surface. Because the microwave signal at L-band
is highly sensitive to the water content in the top few centimetres of soil, these
measurements can be inverted to estimate volumetric soil moisture \citep{kerr1990,smosl2atbd}.
The L2 product delivers this estimate at each node of a fixed global grid (the Discrete Global Grid,
DGG), along with quality flags, retrieval diagnostics, and auxiliary outputs. It is
intended for a wide range of applications including hydrology, meteorology, and climate
research, though known limitations apply over dense forests, frozen/snow-covered soils,
heavily urbanised areas, and regions with strong radio-frequency interference (RFI)
\citep{smosl2atbd}.
\\[1em]
\noindent The L2 product is created through an iterative, physics-based retrieval
algorithm that takes the L1c brightness temperature observations as input and fits them
to a forward radiative-transfer model \citep{kerr2012}. For each DGG node, the algorithm first
pre-processes the L1c views~--- filtering for quality and RFI, computing incidence
angles, and co-locating a range of auxiliary data (land cover, ECMWF meteorological
fields, LAI, soil texture, etc.) on a fine sub-grid. A decision tree then classifies
the scene (vegetated soil, open water, snow, frozen soil, urban, etc.) and selects the
appropriate forward model \citep{kerr2012}. An iterative minimisation (Levenberg--Marquardt) is run to
find the values of soil moisture and vegetation opacity that best fit the observed
multi-angular, dual-polarisation brightness temperatures \citep{smosl2atbd}.
\\[1em]
\noindent The following list contains all data filters applied to \detokenize{$<Run1ContentVars['ConfigVars']['DS1']['name']>$}
before any metrics are computed.

For a detailed description and threshold values for each filter, see the validation run online.

\begin{itemize}
\detokenize{$<'.'.join([f"\\item {d['description']}" for d in Run1ContentVars['ConfigVars']['DS1']['basic_filters_description']])>$}
\end{itemize}

% ---------------------------------------------------------------------------------------------------------------------

\subsection{Reference Data 1 -- ERA5-Land (\detokenize{$<Run2ContentVars['ConfigVars']['DS0']['version']>$}) Surface SM}

\noindent ERA5-Land \citep{munozsabater2021}, produced by ECMWF, is a global reanalysis available from 1950 to present.
It provides surface variables with an increased spatial resolution compared to ERA5 \citep{hersbach2020}.
Soil Moisture in ERA5-Land is available in 1-hour intervals on an approx. 9 km grid and without temporal gaps
at four depth layers (at 0-7 cm depth, 7-28 cm, 28-100 cm, and 100-289 cm).
\\[1em]
Here we use ERA5-Land data in the above-mentioned validation period extracted at hours 0, 6, 12, and 18 of each day at 0–7 cm depth.
Soil temperature information from the same layer is used to mask observations before validation
(when $T < 0$ °C) to eliminate the potential impact of remaining outliers in the satellite data at the transition
from liquid to frozen soil moisture. ERA5-Land data are available as a preliminary release (ERA5-Land T) with a delay
of only a few days.
\\[1em]
The data are then evaluated and finally released with a delay of approximately two months from the current date.
Only the final version of the data is currently used in QA4SM.
\\[1em]
\noindent The top-layer soil moisture field of ERA5-Land is compared to SMOS L2.
The following filters are applied to ERA5-Land before it is temporally matched to SMOS L2 to compute the validation metrics:

\begin{itemize}
\detokenize{$<'.'.join([f"\\item {d['description']}" for d in Run2ContentVars['ConfigVars']['DS0']['basic_filters_description']])>$}
\end{itemize}

% ---------------------------------------------------------------------------------------------------------------------

\subsection{Reference Data 2 -- C3S Satellite SM (COMBINED) (\detokenize{$<Run1ContentVars['ConfigVars']['DS0']['version']>$}) Soil Moisture}

\noindent C3S SM is a merged multi-satellite soil moisture record spanning the period from 1978 to present day. The dataset
harmonizes and combines the observations from currently 19 active (radar) and passive (radiometer) sensors into a single,
consistent long-term record.
\\[1em]
Satellite soil moisture measurements are representative of the first 3-5 centimeters of soil -- depending on the wavelength
used for retrieval, which varies between sensors merged in C3S SM. Data are retrieved at an approximately 25 km spatial
resolution at daily temporal sampling (containing both day- and nighttime retrievals).
\\[1em]
C3S provides regular temporal extensions of the satellite soil moisture time series as Intermediate Climate Data Records (ICDRs).
These usually lag behind present day by 10-20 days.
\\[1em]
The surface soil moisture variable ("sm") of C3S SM is used in the comparison to C3S SM. The following filters are applied to
C3S SM before it is temporally matched to SMOS L2 to compute the validation metrics:

\begin{itemize}
\detokenize{$<'.'.join([f"\\item {d['description']}" for d in Run1ContentVars['ConfigVars']['DS0']['basic_filters_description']])>$}
\end{itemize}
